% !TEX root = ../main.tex

\chapter{Pseudo-holomorphic curves as a derived moduli problem}
\label[chapter]{chap:Pseudo_Holomorphic_Curves_Derived_Moduli}

Let $(\Sigma,j)$ be a closed Riemann surface and $(Z,J)$ an almost complex
manifold. A map $u\colon\Sigma\to Z$ is pseudo-holomorphic when its derivative
is complex linear. This first-order equation provides an application of the
elliptic representability argument of the preceding three chapters.

We first write the equation as a section of its bundle of
complex-antilinear forms and put it into the local coordinates required
by the representability proof. The main calculation is its principal
symbol. Ellipticity then gives a quasi-smooth derived scheme of maps
from a fixed domain, with tangent complex given by the linearized
Cauchy--Riemann equation. Constant maps provide an explicit example
of its deformation and obstruction groups.

Allowing the domain to vary produces a derived moduli stack over the
smooth stack of Riemann surfaces. This relative formulation keeps
deformations of the map distinct from variations and automorphisms of
its domain. We conclude by explaining what further input would be
needed for nodal compactifications and enumerative invariants.

\section{The Cauchy--Riemann equation}
\label[section]{sec:Cauchy_Riemann_Equation_Fourteen}

The condition that a derivative be complex linear is expressed by
the vanishing of its complex-antilinear part. We recall this
decomposition and apply it pointwise to a smooth map.

An almost complex structure on a smooth manifold $Z$ is a vector-bundle
endomorphism
\[
  J\colon TZ\longrightarrow TZ
\]
satisfying $J^2=-\operatorname{id}_{TZ}$. Each tangent space $T_xZ$ thereby
becomes a complex vector space, with multiplication by $i$ given by $J_x$.
The adjective \emph{almost} records that these pointwise complex structures need
not arise from complex coordinates on $Z$.

A Riemann surface can be described similarly as a smooth real surface
$\Sigma$ equipped with an endomorphism $j\colon T\Sigma\to T\Sigma$ satisfying
$j^2=-1$. In real dimension two every almost complex structure is integrable,
but this fact will not be needed. The tangent-bundle formulation puts the
differential equation directly in view.

We begin with the relevant linear algebra. Let $(V,j)$ and $(W,J)$ be complex
vector spaces regarded as real vector spaces with chosen complex structures.
For a real-linear map $L\colon V\to W$, define
\begin{equation}
  \label{eq:Linear_Antilinear_Decomposition_Fourteen}
  L^{1,0}=\frac12(L-JLj),
  \qquad
  L^{0,1}=\frac12(L+JLj).
\end{equation}

\begin{lemma}[Complex-linear decomposition]
  \label[lemma]{lem:Complex_Linear_Decomposition_Fourteen}
  The map $L^{1,0}$ is complex linear, the map $L^{0,1}$ is complex
  antilinear, and
  \[
    L=L^{1,0}+L^{0,1}.
  \]
  Consequently, $L$ is complex linear if and only if $L^{0,1}=0$.
\end{lemma}

\begin{proof}
  Using $j^2=J^2=-1$ gives
  \[
    L^{1,0}j=\tfrac12(Lj+JL)=JL^{1,0},\qquad
    L^{0,1}j=\tfrac12(Lj-JL)=-JL^{0,1}.
  \]
  The sum of the two parts is $L$. Thus $L^{0,1}=0$ implies that
  $L$ is complex linear; conversely, $Lj=JL$ gives $JLj=-L$ and
  hence $L^{0,1}=0$.
\end{proof}

For $p\in\Sigma$ and $x\in Z$, write
\[
  \Hom^{0,1}_{j,J}(T_p\Sigma,T_xZ)
\]
for the real-linear maps $A\colon T_p\Sigma\to T_xZ$ satisfying
$A\circ j=-J\circ A$. These vector spaces form a vector bundle over
$\Sigma\times Z$. Restriction along the graph of a map $u$ gives the bundle
of complex-antilinear one-forms with values in $u^*TZ$, whose sections are
denoted
\[
  \Omega^{0,1}(\Sigma,u^*TZ)
  =
  \Gamma\bigl(
    \Hom^{0,1}_{j,J}(T\Sigma,u^*TZ)
  \bigr).
\]

\begin{definition}[Nonlinear Cauchy--Riemann operator]
  \label[definition]{def:Nonlinear_Cauchy_Riemann_Operator_Fourteen}
  For a smooth map $u\colon\Sigma\to Z$, define
  \begin{equation}
    \label{eq:Nonlinear_Cauchy_Riemann_Operator_Fourteen}
    \bar\partial_Ju
    =
    \frac12\bigl(du+J\circ du\circ j\bigr)
    \in\Omega^{0,1}(\Sigma,u^*TZ).
  \end{equation}
  The map $u$ is \emph{$J$-pseudo-holomorphic}, or
  \emph{$J$-holomorphic}, if $\bar\partial_Ju=0$.
\end{definition}

By \Cref{lem:Complex_Linear_Decomposition_Fourteen}, the equation is
equivalent to
\[
  du\circ j=J\circ du.
\]
When $Z=\mathbb C^n$ with its standard complex structure, this is the ordinary
Cauchy--Riemann equation in local coordinates.

The dependence on the value $u(p)$ makes this a nonlinear equation in
general: varying $u$ changes the endomorphism $J_{u(p)}$ as well as $du$.

The equation is invariant under reparametrization. If
$\varphi\colon(\Sigma',j')\to(\Sigma,j)$ is biholomorphic, then
\[
  \bar\partial_J(u\circ\varphi)=0
\]
whenever $\bar\partial_Ju=0$. Indeed,
$d\varphi\circ j'=j\circ d\varphi$, and substitution in
\Cref{eq:Nonlinear_Cauchy_Riemann_Operator_Fourteen} proves the claim. This
elementary invariance is the first indication that a moduli object for varying
domains must retain automorphisms.

Steffens formulates the application for a symplectic manifold $(Z,\omega)$
equipped with an $\omega$-tame almost complex structure, meaning
\[
  \omega(v,Jv)>0
\]
for every nonzero tangent vector $v$. The symplectic form is fundamental for
energy estimates and compactness, but the equation, the symbol calculation,
and the local representability argument below use only $J$. We will state the
representability theorem first for an almost complex target and return to the
role of $\omega$ at the end.

\section{The fixed-domain solution stack}
\label[section]{sec:Cauchy_Riemann_Local_Coordinates}

The output of the Cauchy--Riemann operator lies in a bundle depending
on the map. First jets define the global section and its derived zero
locus. Near any chosen map, trivializing the equation bundle gives a
single operator between sections of fixed vector bundles.

The equation depends only on the value and first derivative of $u$. A first
jet at $p\in\Sigma$ can be represented by a pair
\[
  (x,L),
  \qquad
  x\in Z,
  \quad
  L\colon T_p\Sigma\longrightarrow T_xZ.
\]
The complex-antilinear projection defines a smooth map on first jets,
\begin{equation}
  \label{eq:Cauchy_Riemann_Jet_Map_Fourteen}
  (x,L)
  \longmapsto
  \left(x,\frac12(L+J_xLj_p)\right).
\end{equation}
Thus pseudo-holomorphicity is a first-order differential equation in the sense
of \Cref{chap:Section_Stacks_Jets_Solution_Stacks}.

For the moment keep the domain fixed. The smooth mapping object
\[
  \mathcal B=C^\infty(\Sigma,Z)
\]
has tangent space
\[
  T_u\mathcal B=\Gamma(\Sigma,u^*TZ)
\]
at $u$. Over $\mathcal B$ there is a bundle $\mathcal F\to\mathcal B$ whose
fiber is
\[
  \mathcal F_u=\Omega^{0,1}(\Sigma,u^*TZ).
\]
The nonlinear Cauchy--Riemann operator is a section
\[
  \bar\partial_J\colon\mathcal B\longrightarrow\mathcal F.
\]
To take its derived zero locus, regard these as derived section stacks:
\[
 \mathcal B=\Sec_\Sigma(\Sigma\times Z),\qquad
 \mathcal F=\Sec_\Sigma(H^{0,1}),
\]
where $H^{0,1}\to\Sigma\times Z$ is the bundle defined above and the second
section stack uses its composite map to $\Sigma$.
The comparison in \Cref{chap:Elliptic_Representability_Smooth_Derived_Bridge} identifies their smooth-stack restrictions
with the indicated convenient manifolds.
The derived solution object is
$\mathcal B\times_{\mathcal F}^{\mathrm{der}}\mathcal B$, using
$\bar\partial_J$ and the zero section.

Fix a map $u$. Exponential coordinates identify nearby maps with a
neighbourhood $Q$ of zero in $\Gamma(\Sigma;F)$, where $F=u^*TZ$.
Choose a connection preserving $J$. Parallel transport along the short
exponential paths gives complex-linear isomorphisms
$\tau_v\colon u^*TZ\to(\exp_u v)^*TZ$.
These identify the nearby bundles of complex-antilinear forms with
\[
 E=\Hom^{0,1}_{j,J}(T\Sigma,u^*TZ).
\]
The transformed equation is
\begin{equation}\label{eq:Local_Cauchy_Riemann_Operator_Fourteen}
 P(v)=\tau_v^{-1}\bar\partial_J(\exp_u v)
 \colon Q\longrightarrow\Gamma(\Sigma;E).
\end{equation}
All coordinate changes depend pointwise on $v$. Thus $P$ is still a
first-order differential operator between sections of fixed vector bundles.
The local section-stack comparison identifies its derived zero fibre with
an open part of the original solution stack.

This is the form needed for the argument of \Cref{chap:Elliptic_Equations_Kuranishi_Models,chap:Elliptic_Representability_Smooth_Derived_Bridge,chap:Elliptic_Representability_Derived_Charts}.
The first-jet presentation in Steffens's Construction~4.0.1 requires a
correction to its output bundle and a local reduction before applying the
fixed-order ellipticity criterion. The correction and the local reduction are discussed in
\Cref{sec:Pseudoholomorphic_Summary_Supplements_Fourteen}.

\section{Linearization and ellipticity}
\label[section]{sec:Cauchy_Riemann_Linearization_Fourteen}

We differentiate at a solution and calculate the principal symbol of
the resulting operator. The symbol is invertible on every nonzero
covector, so the elliptic theory of the preceding chapter applies.
The index formula and the constant-map example describe the resulting
finite-dimensional linear data.

Let $u$ be pseudo-holomorphic. Because $\bar\partial_Ju=0$, the vertical
derivative of the section at $u$ is intrinsically defined:
\[
  D_u\colon
  \Gamma(\Sigma,u^*TZ)
  \longrightarrow
  \Omega^{0,1}(\Sigma,u^*TZ).
\]
Choose a torsion-free connection $\nabla$ on $TZ$ in order to write it.

\begin{proposition}[Linearized Cauchy--Riemann operator]
  \label[proposition]{prop:Linearized_Cauchy_Riemann_Operator_Fourteen}
  For $\xi\in\Gamma(\Sigma,u^*TZ)$, the linearization is
  \begin{equation}
    \label{eq:Linearized_Cauchy_Riemann_Operator_Fourteen}
    D_u\xi
    =
    \frac12\bigl(
      \nabla\xi+J\circ\nabla\xi\circ j
    \bigr)
    +
    \frac12(\nabla_\xi J)\circ du\circ j.
  \end{equation}
  In particular, $D_u$ is a real Cauchy--Riemann type operator: its
  first-order part is a Cauchy--Riemann operator, while the second summand is
  zeroth order.
\end{proposition}

\begin{proof}
  Choose a smooth family $u_t$ with $u_0=u$ and
  $\frac{d}{dt}|_{t=0}u_t=\xi$. Covariant differentiation gives
  \[
    \left.\frac{D}{dt}\right|_{t=0}du_t=\nabla\xi.
  \]
  Differentiating the target almost complex structure gives
  \[
    \left.\frac{D}{dt}\right|_{t=0}J(u_t)=\nabla_\xi J.
  \]
  The complex structure $j$ on the fixed domain does not vary. Applying the
  product rule to
  \Cref{eq:Nonlinear_Cauchy_Riemann_Operator_Fourteen} proves
  \Cref{eq:Linearized_Cauchy_Riemann_Operator_Fourteen}.

  The connection is used only to compare the bundles $u_t^*TZ$ for varying
  $t$. The vertical derivative of a bundle section at a zero is independent
  of this choice, so $D_u$ is intrinsic.
\end{proof}

The first summand differentiates $\xi$. The term involving $\nabla_\xi J$
contains no derivative of $\xi$, and hence contributes nothing to the
principal symbol. Consequently the principal symbol is the same as that of the first summand.

For a first-order differential operator, the principal symbol is obtained by
retaining only the coefficient of the first derivative. Fix $p\in\Sigma$,
$0\neq\zeta\in T_p^*\Sigma$, and $v\in T_{u(p)}Z$.

\begin{theorem}[Ellipticity of the Cauchy--Riemann operator]
  \label[theorem]{thm:Cauchy_Riemann_Ellipticity_Fourteen}
  The principal symbol of $D_u$ is
  \begin{equation}
    \label{eq:Cauchy_Riemann_Principal_Symbol_Fourteen}
    \sigma_\zeta(D_u)(v)(Y)
    =
    \frac12\bigl(
      \zeta(Y)v+\zeta(jY)Jv
    \bigr).
  \end{equation}
  For every nonzero $\zeta$, this is an isomorphism
  \[
    T_{u(p)}Z
    \iso
    \Hom^{0,1}_{j,J}(T_p\Sigma,T_{u(p)}Z)
  \]
  of real vector spaces. Consequently, $D_u$ is elliptic.
\end{theorem}

\begin{proof}
  The first-order part of $D_u$ is
  \[
    \xi\longmapsto
    \frac12(\nabla\xi+J\circ\nabla\xi\circ j).
  \]
  Replacing $\nabla_Y\xi$ by $\zeta(Y)v$ gives
  \Cref{eq:Cauchy_Riemann_Principal_Symbol_Fourteen}.

  First observe that the resulting map is complex antilinear in $Y$. Indeed,
  \begin{align*}
    \sigma_\zeta(D_u)(v)(jY)
    &=\frac12\bigl(
       \zeta(jY)v-\zeta(Y)Jv
      \bigr)\\
    &=-J\sigma_\zeta(D_u)(v)(Y).
  \end{align*}

  Choose a $j$-compatible inner product on $T_p\Sigma$ and let $Y$ be the
  metric dual of $\zeta$. Then
  \[
    \zeta(Y)>0,
    \qquad
    \zeta(jY)=0.
  \]
  If $\sigma_\zeta(D_u)(v)=0$, evaluation at $Y$ gives
  \[
    0=\frac12\zeta(Y)v,
  \]
  and therefore $v=0$. The symbol is injective. Since $T_p\Sigma$ has complex
  dimension one, a complex-antilinear map out of $T_p\Sigma$ is determined by
  its value on one nonzero real direction. Hence the source and target of the
  symbol have the same real dimension, so the symbol is an isomorphism.
\end{proof}

This proof is pointwise. It uses neither a symplectic form on $Z$ nor
integrability of $J$. It also works uniformly when the domain and the map vary
in smooth families.

Since $\Sigma$ is closed, after completing in Sobolev norms,
the elliptic Fredholm theorem applies.

\begin{corollary}[Fredholm deformation operator]
  \label[corollary]{cor:Cauchy_Riemann_Fredholm_Fourteen}
  For sufficiently large $k$, the extension
  \[
    D_u\colon
    H^{k+1}(\Sigma,u^*TZ)
    \longrightarrow
    H^k\Omega^{0,1}(\Sigma,u^*TZ)
  \]
  is Fredholm. Its kernel and cokernel are finite-dimensional and are
  represented by smooth sections.
\end{corollary}

\begin{proof}
  Fredholmness follows from
  \Cref{thm:Cauchy_Riemann_Ellipticity_Fourteen} and the elliptic Fredholm
  theorem used in
  \Cref{thm:Elliptic_Fredholm_Regularity}.
  Elliptic regularity identifies the Sobolev kernel and cokernel with their
  smooth counterparts.
\end{proof}

If $Z$ has complex dimension $n$, Riemann--Roch gives the real index
\begin{equation}
  \label{eq:Cauchy_Riemann_Index_Fourteen}
  \operatorname{ind}_{\mathbb R}(D_u)
  =
  n\chi(\Sigma)
  +2\left\langle c_1(u^*TZ),[\Sigma]\right\rangle.
\end{equation}
This is the Riemann--Roch formula
\cite[Theorem~3.22]{Wendl_Lectures_Holomorphic_Curves}.
It is the virtual dimension of the moduli problem
relative to the fixed domain. For an unparametrized problem with varying
domains, the deformation theory of the domain must also be included.

\begin{example}[Constant maps]
Fix $x\in Z$ and let $u_x\colon\Sigma\to Z$ be the constant map with value
$x$. Since $du_x=0$, it is pseudo-holomorphic, and the zeroth-order term in
\Cref{eq:Linearized_Cauchy_Riemann_Operator_Fourteen} vanishes. Thus
$D_{u_x}$ is the Dolbeault operator on the trivial complex vector bundle with
fiber $T_xZ$:
\[
  \bar\partial\colon
  \Omega^0(\Sigma)\otimes_{\mathbb C}T_xZ
  \longrightarrow
  \Omega^{0,1}(\Sigma)\otimes_{\mathbb C}T_xZ.
\]
For connected $\Sigma$ this gives
\begin{equation}
  \label{eq:Constant_Map_Deformation_Obstruction_Fourteen}
  \ker D_{u_x}\cong T_xZ,
  \qquad
  \operatorname{coker}D_{u_x}
  \cong
  H^{0,1}(\Sigma)\otimes_{\mathbb C}T_xZ.
\end{equation}

Indeed, holomorphic functions on a closed connected Riemann surface are
constant, giving the kernel. Dolbeault cohomology is, by definition, the
cokernel of the scalar $\bar\partial$ operator on a surface, giving the
second formula. If $\Sigma$ has genus $g$ and $\dim_{\mathbb C}Z=n$, these
groups have real dimensions $2n$ and $2ng$. The index is $2n(1-g)$.

For a sphere, the cokernel vanishes. For a positive-genus domain,
constant maps have nonzero obstruction space, although the visible family of
constant maps is the ordinary manifold $Z$.
The calculation concerns the full mapping problem at a constant map;
it does not identify the entire moduli object with the constant-map locus.
\end{example}

These kernel and cokernel calculations give the linear data.
The Kuranishi construction identifies their two-term complex with
the intrinsic derived tangent complex, as we now explain.

\section{Deformation theory with the domain fixed}
\label[section]{sec:Pseudo_Holomorphic_Deformation_Theory_Fourteen}

The local Kuranishi construction turns the elliptic deformation
complex into an intrinsic tangent complex. We use this identification
to interpret regularity, obstruction groups, and virtual dimension
while keeping the domain fixed.

Let
\[
  \mathcal M_{(\Sigma,j)}(Z,J)
\]
denote the derived solution stack for the fixed domain $(\Sigma,j)$. On
ordinary test manifolds, a $T$-point is a smooth map
\[
  u\colon T\times\Sigma\longrightarrow Z
\]
whose restriction to every slice is pseudo-holomorphic. On a derived test
object, the equation is imposed by a derived zero locus and hence retains its
infinitesimal intersection data.

The analytic deformation complex is the two-term complex with differential
$D_u$. Its identification with a derived tangent complex follows from the
finite-dimensional Kuranishi reduction.

\begin{theorem}[Fixed-domain tangent complex]
  \label[theorem]{thm:Fixed_Domain_Tangent_Complex_Fourteen}
  At a pseudo-holomorphic map $u\colon(\Sigma,j)\to(Z,J)$, the tangent complex
  of the fixed-domain derived solution stack is
  \begin{equation}
    \label{eq:Fixed_Domain_Tangent_Complex_Fourteen}
    \mathbb T_u\mathcal M_{(\Sigma,j)}(Z,J)
    \simeq
    \left[
      \Gamma(\Sigma,u^*TZ)
      \xrightarrow{D_u}
      \Omega^{0,1}(\Sigma,u^*TZ)
    \right]
  \end{equation}
  in homological degrees $0$ and $-1$.
\end{theorem}

\begin{proof}
  Use the local operator of
  \Cref{sec:Cauchy_Riemann_Local_Coordinates}.
  Its derivative at the chosen zero is $D_u$, which is elliptic.
  The construction of \Cref{chap:Elliptic_Representability_Derived_Charts} gives a finite-dimensional chart
  $\kappa\colon Z_U\to V$ and a quasi-isomorphism
  \[
    [T_uZ_U\longrightarrow V]\longrightarrow
    [\Gamma(\Sigma,u^*TZ)\xrightarrow{D_u}
      \Omega^{0,1}(\Sigma,u^*TZ)].
  \]
  The left side is the tangent complex of the finite-dimensional derived
  zero locus, by the calculation in \Cref{chap:Stable_Linearization_And_Tangent_Complexes}. This proves the formula
  without applying a finite-dimensional cotangent theorem directly to
  the infinite-dimensional mapping object.
\end{proof}

Consequently,
\begin{align*}
  H_0\mathbb T_u
  &\cong\ker D_u,\\
  H_{-1}\mathbb T_u
  &\cong\operatorname{coker}D_u.
\end{align*}
The first group is the space of infinitesimal deformations of the map with the
domain fixed. The second is its obstruction space. Both are finite-dimensional
by \Cref{cor:Cauchy_Riemann_Fredholm_Fourteen}.

\begin{definition}[Regular pseudo-holomorphic map]
  \label[definition]{def:Regular_Pseudo_Holomorphic_Map_Fourteen}
  A pseudo-holomorphic map $u$ is \emph{regular} if $D_u$ is surjective.
\end{definition}

At a regular map the obstruction group vanishes. The Banach implicit function
theorem then identifies the nearby ordinary solution space with a smooth
manifold of dimension $\operatorname{ind}_{\mathbb R}(D_u)$. Derivedly, the
tangent complex is concentrated in degree $0$.

If the cokernel is nonzero, it receives obstructions to extending an
infinitesimal deformation to an actual family. A particular obstruction
class can still vanish. The Kuranishi map describes these nonlinear
conditions; its derivative alone need not determine them.

The virtual dimension at $u$ is
\[
  \dim H_0\mathbb T_u-\dim H_{-1}\mathbb T_u
  =\operatorname{ind}_{\mathbb R}(D_u),
\]
which is given by \Cref{eq:Cauchy_Riemann_Index_Fourteen}.

\section{Varying domains and relative representability}
\label[section]{sec:Pseudo_Holomorphic_Families_Representability_Fourteen}

We now retain the complex structure and automorphisms of the domain
as part of the moduli problem. The stack of compact Riemann surfaces
supplies a universal proper family. The relative Cauchy--Riemann
section defines the solution stack over this base, and the local
elliptic argument proves its relative representability.

The fixed-domain problem ignores reparametrizations and variations of the
complex structure. To retain both, let
\[
  S=\operatorname{Surf}_{\mathbb C}
\]
denote the smooth stack of compact Riemann surfaces. For a test manifold $T$,
an object of $S(T)$ is a proper submersion
\[
  \pi\colon C\longrightarrow T
\]
with compact surface fibers, together with a smooth complex structure on the
vertical tangent bundle $T(C/T)$. Morphisms are Cartesian diagrams whose maps
on fibers are biholomorphic.

The stack retains smooth dependence on parameters and the automorphisms of
each surface. For example, the Riemann sphere has automorphism group
$\operatorname{PSL}_2(\mathbb C)$, which appears as isotropy at its point of $S$.

The stack carries a universal family
\begin{equation}
  \label{eq:Universal_Riemann_Surface_Family_Fourteen}
  M=\widetilde{\operatorname{Surf}}_{\mathbb C}
  \longrightarrow
  S=\operatorname{Surf}_{\mathbb C}.
\end{equation}
Its fiber over $(\Sigma,j)$ is the surface $\Sigma$ itself. In particular,
$M\to S$ is a proper family of smooth compact manifolds, which is exactly the
domain hypothesis in the elliptic representability theorem.

Kodaira--Spencer theory implies that $S$ is a smooth Artin
$\Cinfty$-stack. We quote this result, as Steffens does in
\cite[Remark~4.0.3]{Steffens_Representability_Elliptic_Moduli}. Its role will
be to turn relative representability over $S$ into an absolute derived Artin
statement.

The Artin terminology allows smooth submersive atlases, with representable
overlaps, in place of the open atlases for schemes. Such a stack can have
positive-dimensional groups of automorphisms. In the derived version the
atlas objects are derived schemes. Pulling a smooth atlas of $S$ back to a
relatively representable moduli stack gives an atlas of this kind.
Here a derived Artin stack is called quasi-smooth when it admits a smooth
atlas by quasi-smooth derived schemes.

Put $Y=Z\times M$ over $M$. A section of $Y$ over a family $C\to T$ is a
smooth family of maps $C\to Z$. The corresponding section stack
\[
  \mathcal B
  =\Sec_{M/S}(Y)
\]
parametrizes the unknown maps.

Over $Y$ there is a finite-rank vector bundle $H^{0,1}\to Y$ whose fiber at
$(c,x)$ is
\[
  \Hom^{0,1}_{j_c,J_x}
  \bigl(T_c(M/S),T_xZ\bigr).
\]
Pulling this bundle back along a map-section and taking relative sections
produces a bundle $\mathcal F\to\mathcal B$. Its fiber at a family $u$ is the
family of spaces
\[
  \Omega^{0,1}(C_t,u_t^*TZ).
\]

The complex-antilinear projection of the relative first jet defines a section
\[
  \operatorname{CR}\colon\mathcal B\longrightarrow\mathcal F,
  \qquad
  u\longmapsto\bar\partial_Ju.
\]
The construction uses only the vertical tangent bundle, the fiberwise complex
structure, and $J$. It is therefore natural under Cartesian base change of
families of surfaces.

\begin{definition}[Derived moduli stack of pseudo-holomorphic maps]
  \label[definition]{def:Derived_Moduli_Pseudoholomorphic_Maps_Fourteen}
  The \emph{derived moduli stack of pseudo-holomorphic maps}, denoted
  $\mathcal M(Z,J)$, is the derived zero locus of the Cauchy--Riemann section
  over $\mathcal B$.
\end{definition}

It comes with a natural morphism
\[
  \mathcal M(Z,J)
  \longrightarrow
  \operatorname{Surf}_{\mathbb C}.
\]
Over an ordinary test manifold $T\to S$ classifying $C\to T$, its ordinary
points are precisely smooth maps $u\colon C\to Z$ whose restrictions to the
fibers are pseudo-holomorphic. The derived zero locus retains the higher
intersection data of this equation.

\begin{theorem}[Pseudo-holomorphic representability]
  \label[theorem]{thm:Pseudoholomorphic_Representability_Fourteen}
  Let $(Z,J)$ be an almost complex manifold. The morphism
  \begin{equation}
    \label{eq:Pseudoholomorphic_Representability_Morphism_Fourteen}
    \mathcal M(Z,J)
    \longrightarrow
    \operatorname{Surf}_{\mathbb C}
  \end{equation}
  is representable by quasi-smooth derived $\Cinfty$-schemes locally of finite
  presentation. Consequently, $\mathcal M(Z,J)$ is a quasi-smooth derived
  Artin $\Cinfty$-stack.
\end{theorem}

\begin{proof}
  Representability is local on the target. Pull back along a map
  $T\to\operatorname{Surf}_{\mathbb C}$ from a manifold. It classifies a
  smooth proper family of Riemann surfaces $C\to T$. Locally on $T$, Ehresmann
  triviality identifies the underlying family with
  $T'\times\Sigma\to T'$, although the fiberwise complex structure may still
  vary with the parameter.

  Fix a solution $u$ in this family. The local section-stack construction from
  \Cref{sec:Cauchy_Riemann_Local_Coordinates} identifies a neighborhood of $u$ with
  $T'\times Q$, where $Q$ is an open subset of a smooth section space. After
  trivializing the nearby equation bundles, the family Cauchy--Riemann section
  becomes a smooth family of nonlinear first-order operators
  \begin{equation}
    \label{eq:Family_Local_Cauchy_Riemann_Operator_Fourteen}
    P\colon
    T'\times Q
    \longrightarrow
    T'\times\Gamma(\Sigma;E).
  \end{equation}
  Its vertical linearization at $u$ is $D_u$, which is elliptic by
  \Cref{thm:Cauchy_Riemann_Ellipticity_Fourteen}.

  The local Kuranishi construction of \Cref{chap:Elliptic_Representability_Smooth_Derived_Bridge,chap:Elliptic_Representability_Derived_Charts} applies to
  this operator. It produces a quasi-smooth derived
  $\Cinfty$-scheme locally of finite presentation representing a neighborhood
  of $u$ in the solution stack. These neighborhoods cover the pullback over
  $T$, and the open-atlas criterion glues them. This proves relative
  representability.

  Finally, $\operatorname{Surf}_{\mathbb C}$ is a smooth Artin
  $\Cinfty$-stack. Relative quasi-smooth representability over this base makes
  the total solution stack a quasi-smooth derived Artin $\Cinfty$-stack.
\end{proof}

Steffens states this conclusion in the symplectic tame setting in
\cite[Construction~4.0.1 and Remark~4.0.3]{Steffens_Representability_Elliptic_Moduli}.
The proof above uses the local zero-section form of the representability
argument, for the reason explained in
\Cref{warn:Steffens_Five_Tuple_Fourteen}. It also shows that the local
representability statement depends only on $J$.

In particular, given a manifold $T$ and a
specified family $C\to T$ of Riemann surfaces, the pullback
\[
  T\times_{\operatorname{Surf}_{\mathbb C}}\mathcal M(Z,J)
\]
is represented by a quasi-smooth derived $\Cinfty$-scheme locally of finite
presentation over $T$. Its relative tangent complex at a map $u$ is
\begin{equation}
  \label{eq:Relative_Tangent_Complex_Final_Fourteen}
  \mathbb T_u\bigl(
    \mathcal M(Z,J)/\operatorname{Surf}_{\mathbb C}
  \bigr)
  \simeq
  \left[
    \Gamma(\Sigma,u^*TZ)
    \xrightarrow{D_u}
    \Omega^{0,1}(\Sigma,u^*TZ)
  \right].
\end{equation}
The absolute tangent complex also contains deformations and infinitesimal
automorphisms of the domain. The corresponding fiber sequence is recorded in
\Cref{eq:Absolute_Tangent_Fiber_Sequence_Fourteen}.

\section{Compactification and enumerative questions}
\label[section]{sec:Pseudoholomorphic_Boundary_Fourteen}

Properness of the domain family ensures compact fibres on which elliptic
analysis applies. Compactness of the solution object is another issue.
Even with compact target, sequences of maps can develop bubbles, and smooth
domains can degenerate to nodal curves. These limits require an enlarged
moduli problem, typically stable nodal maps with fixed genus, marked points,
and homology class.

Representing that compactification requires analysis near its nodal objects,
including gluing. An enumerative application then needs a properness theorem
and appropriate orientations and virtual classes. \Cref{chap:Derived_Bordism_And_Fundamental_Classes} supplies a
bordism class for a compact quasi-smooth derived manifold. The total
moduli stack here may have isotropy, so a virtual-class construction for
stacks requires additional theory.

When $J$ is tamed by a symplectic form $\omega$, the area
$\int_\Sigma u^*\omega$ is positive for a nonconstant pseudo-holomorphic map
from a connected domain, and depends only on its homology class.
With the associated $J$-invariant metric, it equals the energy of the map.
Thus fixing the class gives an
energy bound, an essential input to compactness. The bound allows bubbling
and explains why nodal maps enter the compactification.

\section{Further reading}
\label[section]{sec:Pseudoholomorphic_Summary_Supplements_Fourteen}

The following discussion compares the local equation used in the proof
with the symmetric presentation in the source, records the absolute
tangent complex when the domain varies, and describes further choices
of geometric data.

\begin{warning}[The printed five-tuple]
  \label[warning]{warn:Steffens_Five_Tuple_Fourteen}
  In Steffens's Construction~4.0.1, the common output of the displayed
  Cauchy--Riemann and zero maps is the full relative first-jet bundle. Taken
  literally, its linearized matching complex has complementary
  complex-linear jet directions in degree minus one, in addition to the
  Cauchy--Riemann Fredholm complex.

  Replacing the common output by the bundle of complex-antilinear maps
  removes those extra obstruction directions. The resulting symmetric matching operator still
  has a zeroth-order component, equating the underlying map values, and a
  first-order Cauchy--Riemann component. It is therefore not elliptic in the
  fixed-order sense of \cite[Definition~3.4.2]{Steffens_Representability_Elliptic_Moduli}.
  The general theorem cannot be invoked verbatim from that five-tuple.
\end{warning}

The derived intersection of the Cauchy--Riemann section with the zero section
can be displayed as
\begin{equation}
  \label{eq:Symmetric_Cauchy_Riemann_Pullback_Fourteen}
  \begin{tikzcd}[column sep=large]
    \mathcal M(Z,J)
      \rar \dar \drar[pullback]
      & \mathcal B \dar["\operatorname{CR}"] \\
    \mathcal B
      \rar["0"']
      & \mathcal F.
  \end{tikzcd}
\end{equation}
At a solution $u$, put
\[
  A=\Gamma(\Sigma,u^*TZ),
  \qquad
  B=\Omega^{0,1}(\Sigma,u^*TZ).
\]
Linearizing this pullback presentation gives
\begin{equation}
  \label{eq:Symmetric_Tangent_Complex_Fourteen}
  \left[
    A\oplus A
    \xrightarrow{\delta}
    A\oplus B
  \right],
  \qquad
  \delta(\xi_1,\xi_2)
  =
  (\xi_1-\xi_2,D_u\xi_1).
\end{equation}
The first component differentiates equality of the underlying maps, and the
second differentiates the Cauchy--Riemann equation.

\begin{lemma}[Cancellation of the repeated field]
  \label[lemma]{lem:Repeated_Field_Cancellation_Fourteen}
  There is a chain isomorphism
  \[
    [A\oplus A\xrightarrow{\delta}A\oplus B]
    \iso
    [A\xrightarrow{D_u}B]
    \oplus
    [A\xrightarrow{\operatorname{id}}A].
  \]
  Consequently,
  \Cref{eq:Symmetric_Tangent_Complex_Fourteen} is quasi-isomorphic to the
  Cauchy--Riemann deformation complex.
\end{lemma}

\begin{proof}
  On the source use the coordinates
  \[
    (\xi_1,\xi_2)
    \longmapsto
    (\xi_1,\xi_1-\xi_2).
  \]
  Reorder the target $A\oplus B$ as $B\oplus A$. In these coordinates the
  differential is
  \[
    (\xi,r)\longmapsto(D_u\xi,r),
  \]
  which is the displayed direct sum. The identity complex is contractible.
\end{proof}

The cancellation recovers the Cauchy--Riemann deformation complex.
Its identification with the intrinsic tangent complex uses the local
Kuranishi reduction proved above. The first-order symbol of the displayed
matching operator is still $(\xi_1,\xi_2)\mapsto(0,\sigma(D_u)\xi_1)$,
which is not invertible. Passing to the single local equation removes the
repeated map variable before applying elliptic analysis.

The complex in \Cref{eq:Relative_Tangent_Complex_Final_Fourteen} holds the
domain and its complex structure fixed. The absolute tangent complex also
contains the deformation theory of $(\Sigma,j)$. It fits into the relative
tangent fiber sequence
\begin{equation}
  \label{eq:Absolute_Tangent_Fiber_Sequence_Fourteen}
  \mathbb T_u\bigl(
    \mathcal M(Z,J)/\operatorname{Surf}_{\mathbb C}
  \bigr)
  \longrightarrow
  \mathbb T_u\mathcal M(Z,J)
  \longrightarrow
  \mathbb T_{(\Sigma,j)}\operatorname{Surf}_{\mathbb C}.
\end{equation}
The last term records deformations and infinitesimal automorphisms of the
domain. For a smooth Riemann surface, these are measured by
$H^1(\Sigma,T\Sigma)$ and $H^0(\Sigma,T\Sigma)$, respectively,
placed in tangent degrees $0$ and $1$. Thus passing from the relative
complex to the absolute one includes both changes of complex structure
and reparametrizations.
In this example the absolute tangent complex has amplitude in $[-1,1]$;
the degree $1$ term accounts for infinitesimal automorphisms. This is
compatible with quasi-smoothness of the Artin stack as defined through
its smooth atlases.

This distinction prevents an ambiguity in expected-dimension formulas. The
index in \Cref{eq:Cauchy_Riemann_Index_Fourteen} is the relative virtual
dimension for a fixed domain. An unparametrized moduli problem with varying
domains must also account for
$\mathbb T_{(\Sigma,j)}\operatorname{Surf}_{\mathbb C}$.

For connected genus-$g$ domains and no marked points, Riemann--Roch for
$T\Sigma$ gives the real virtual dimension $6g-6$ of the domain stack.
Adding this to the relative index gives
\[
 \operatorname{vdim}_{\R}\mathcal M(Z,J)
 =(n-3)(2-2g)+2\langle c_1(u^*TZ),[\Sigma]\rangle.
\]
In genus zero the contribution $-6$ accounts for the six-dimensional
reparametrization group. This is an expected dimension for the stack,
including its infinitesimal automorphisms.

Several variations fit the same representability argument.
\begin{enumerate}
  \item One may let $J$ vary in a smooth family by enlarging the base stack.
  \item One may allow a smooth family of almost complex target manifolds.
  \item One may add marked points by adding sections of the universal curve to
    the objects classified by the domain stack.
  \item One may formulate equivariant problems over quotient stacks, thereby
    retaining isotropy rather than passing to a naive orbit set.
\end{enumerate}
The first three variants are noted in
\cite[Remark~4.0.2]{Steffens_Representability_Elliptic_Moduli}. All still
concern smooth compact domains.

\begin{exercise}
At a constant map from a genus-$g$ surface to a complex $n$-fold, compute
the dimensions of the relative deformation and obstruction spaces.
Check the Riemann--Roch formula.
\end{exercise}

\begin{exercise}
Write the first-order principal symbol of the symmetric matching operator
in \Cref{eq:Symmetric_Tangent_Complex_Fourteen}.
Show why it fails to be invertible even though its complex is
quasi-isomorphic to the elliptic Cauchy--Riemann complex.
\end{exercise}

The primary derived-geometric source is
\citet{Steffens_Representability_Elliptic_Moduli}.
The construction in
\cite[Construction~4.0.1]{Steffens_Representability_Elliptic_Moduli}
defines the stack of smooth compact
Riemann surfaces, its universal family, the jet-level Cauchy--Riemann map, and
the solution stack. Varying data, marked points, and the derived Artin
conclusion are discussed in
\cite[Remarks~4.0.2--4.0.3]{Steffens_Representability_Elliptic_Moduli}.
The source correction needed for the printed five-tuple is recorded in
\Cref{warn:Steffens_Five_Tuple_Fourteen}.

For the analytic background, Wendl's
\emph{Lectures on Holomorphic Curves in Symplectic and Contact Geometry}
develops almost complex structures and pseudo-holomorphic maps in Section~2.1,
real Cauchy--Riemann type operators and their linearization in
Sections~2.3--2.4, regularity in Section~2.5, and Fredholm and index theory in
Sections~3.3--3.4
\cite{Wendl_Lectures_Holomorphic_Curves}. These sections support the
linearization, symbol, Fredholm, and Riemann--Roch statements used here.
